% Part: sets-functions-relations
% Chapter: infinite
% Section: dedekinds-proof
%
\documentclass[../../../include/open-logic-section]{subfiles}

\begin{document}

\olfileid[id]{sfr}{infinite}{dedekindsproof}

\olsection[``Bukti'' Dedekind]{``Bukti'' Dedekind tentang
Keberadaan Himpunan Tak Hingga}

Dalam bab ini, kita telah menyajikan pembahasan bilangan asli dalam teori
himpunan, dengan menggunakan aljabar Dedekind. Dalam
\olref[arith][ref]{sec}, kita merenungkan makna filosofis aritmetisasi
analisis (di antara hal-hal lain). Sekarang kita perlu merenungkan makna dari
apa yang telah kita capai di sini.

Sepanjang \olref[sfr][arith][]{chap}, kita menganggap bilangan asli sebagai
sesuatu yang telah tersedia dan menggunakannya untuk membangun bilangan
bulat, rasional, dan real secara eksplisit. Dalam bab ini, kita belum
memberikan konstruksi eksplisit bilangan asli. Kita hanya menunjukkan bahwa,
\emph{jika diberikan sembarang himpunan yang tak hingga Dedekind}, kita dapat
mendefinisikan suatu himpunan yang akan berperilaku persis sebagaimana kita
kehendaki~$\Nat$ berperilaku.

Jelaslah bahwa kita tidak dapat mengaku telah menjawab suatu pertanyaan
metafisis, seperti \emph{objek manakah yang merupakan bilangan asli}. Namun,
hal itu justru baik. Lagi pula, dalam \olref[sfr][arith][ref]{sec}, kita
menekankan bahwa kita akan keliru jika menganggap definisi $\Real$ sebagai
himpunan potongan Dedekind sebagai suatu \emph{penemuan}, alih-alih suatu
penetapan yang memudahkan. Pengamatan yang menentukan ialah bahwa potongan
Dedekind mencontohkan sifat-sifat matematis utama bilangan real. Demikian
pula di sini: pengamatan yang menentukan ialah bahwa \emph{setiap} aljabar
Dedekind mencontohkan sifat-sifat matematis utama bilangan asli. (Bahkan,
Dedekind menegaskan pokok ini dengan membuktikan bahwa semua aljabar Dedekind
\emph{isomorfik} (\citeyear[Theorems
132--3]{Dedekind1888}). Maka tidak
mengherankan bahwa banyak ``strukturalis'' kontemporer menyebut Dedekind
sebagai seorang perintis.)

Selain itu, kita telah menunjukkan cara menanamkan teori bilangan asli ke
dalam suatu teori himpunan sederhana yang naif. Teori itu sendiri masih
cukup informal, tetapi (tampaknya) tidak mengandaikan bilangan asli sebagai
sesuatu yang telah tersedia. Jadi, mungkin kita sedang menuju realisasi
proyek ambisius Dedekind sendiri, yang ia jelaskan sebagai berikut:
\begin{quote}
	Dalam ilmu pengetahuan, tidak ada sesuatu yang dapat dibuktikan yang patut
	dipercayai tanpa bukti. Meskipun tuntutan ini tampak masuk akal, saya tidak
	dapat menganggapnya telah dipenuhi bahkan dalam metode-metode paling
	mutakhir untuk meletakkan landasan ilmu yang paling sederhana; yakni,
	bagian logika yang berurusan dengan teori bilangan. Dengan menyebut
	aritmetika (aljabar, analisis) sekadar bagian dari logika, maksud saya ialah
	bahwa saya menganggap konsep bilangan sepenuhnya tidak bergantung pada
	gagasan atau intuisi ruang dan waktu---bahwa saya justru menganggapnya
	sebagai produk langsung hukum-hukum pemikiran murni.
	\citep[preface]{Dedekind1888}
\end{quote}
Gagasan berani Dedekind adalah sebagai berikut. Kita baru saja menunjukkan
cara membangun bilangan asli hanya dengan menggunakan teori himpunan (naif).
Dalam \olref[sfr][arith][]{chap}, kita melihat cara membangun bilangan real
jika bilangan asli dan sedikit teori himpunan telah diberikan. Jadi, mungkin
``aritmetika (aljabar, analisis)'' ternyata ``sekadar bagian dari logika''
(dalam pengertian luas Dedekind atas kata ``logika'').

Itulah gagasannya. Namun, tunggu sejenak. Konstruksi aljabar Dedekind kita
(pengganti bilangan asli yang kita gunakan) bergantung pada syarat bahwa
suatu himpunan yang tak hingga Dedekind ada. (Lihat kembali
\olref[sfr][infinite][dedekind]{thm:DedekindInfiniteAlgebra}.) Jika
keberadaan suatu himpunan yang tak hingga Dedekind tidak dapat ditetapkan
melalui ``logika'' atau ``hukum-hukum pemikiran murni'', proyek tersebut
terhenti.

Jadi, \emph{dapatkah} keberadaan suatu himpunan yang tak hingga Dedekind
ditetapkan oleh ``hukum-hukum pemikiran murni''? Berikut upaya Dedekind:
\begin{quote}
  Ranah pemikiran saya sendiri, yakni totalitas $S$ dari segala hal yang
  dapat menjadi objek pemikiran saya, bersifat tak hingga. Sebab, jika $s$
  menandai suatu elemen dari~$S$, maka pemikiran $s'$ bahwa~$s$ dapat menjadi
  objek pemikiran saya itu sendiri merupakan suatu elemen dari~$S$. Jika kita
  memandangnya sebagai suatu citra $\phi(s)$ dari elemen~$s$, maka
  \dots~$S$ bersifat [tak hingga Dedekind], sebagaimana hendak dibuktikan.
	\citep[\S66]{Dedekind1888}
\end{quote}
Ini merupakan hal yang sangat mencengangkan untuk ditemukan di tengah-tengah
sebuah buku yang sebagian besarnya terdiri atas bukti-bukti matematis yang
sangat ketat. Dua catatan patut dikemukakan.

Pertama: ``bukti'' ini hampir tidak mempunyai sifat yang sekarang akan kita
kenali sebagai sifat ``matematis''. Bukti ini berbicara tentang objek
psikologis (pemikiran), dan bahkan sekadar objek yang \emph{mungkin} ada.

Kedua: sekurang-kurangnya sebagaimana kita telah menyajikan aljabar Dedekind,
``bukti'' ini mempunyai kekurangan teknis yang langsung terlihat. Jika
argumen Dedekind berhasil, argumen itu hanya menetapkan bahwa terdapat hal
yang banyaknya tak hingga (secara khusus, pemikiran yang banyaknya tak
hingga). Namun, Dedekind juga perlu memberi kita alasan untuk memandang~$S$
sebagai satu \emph{himpunan}, dengan !!{element}s yang banyaknya tak hingga,
alih-alih memandang~$S$ sebagai \emph{sejumlah hal} (dalam bentuk jamak).

Fakta bahwa Dedekind tidak melihat kesenjangan di sini mungkin menunjukkan
bahwa penggunaan kata ``totalitas'' olehnya tidak secara tepat sejalan
dengan penggunaan kata ``himpunan'' oleh \emph{kita}.\footnote{Memang, kita
mempunyai alasan-alasan lain untuk berpikir demikian; lihat
\citet[p.~23]{Potter2004}.} Namun, hal ini tidak terlalu mengejutkan. Proyek
yang kita kerjakan dalam dua bab terakhir---suatu ``konstruksi'' bilangan
asli, lalu dari bilangan itu suatu ``konstruksi'' bilangan bulat, real, dan
rasional---seluruhnya dilaksanakan secara naif. Kita begitu saja menggunakan
himpunan ini atau himpunan itu ketika membutuhkannya, tanpa menetapkan banyak
prinsip umum mengenai himpunan mana tepatnya yang ada, dan kapan. Namun, kita
tahu bahwa kita memerlukan \emph{beberapa} prinsip umum; jika tidak, kita
akan terjerumus ke dalam Paradoks Russell.

Kini tibalah waktunya bagi kita untuk meninggalkan kenaifan tersebut.

\end{document}
